\documentclass[12pt,a4paper]{article}
\usepackage[utf8]{inputenc}
\usepackage[T1]{fontenc}
\usepackage[french]{babel}
\usepackage[margin=2.2cm, top=2.5cm, bottom=2.5cm]{geometry}
\usepackage{amsmath,amssymb}
\usepackage{booktabs}
\usepackage{xcolor}
\usepackage{titlesec}
\usepackage{hyperref}
\usepackage{lmodern}
\usepackage{fancyhdr}
\usepackage{tcolorbox}
\usepackage{enumitem}
\usepackage{array}
\usepackage{tabularx}
\tcbuselibrary{skins,breakable}
\setlength{\parskip}{5pt}
\setlength{\parindent}{0pt}

% Colors
\definecolor{primary}{RGB}{22,60,110}
\definecolor{accent}{RGB}{180,40,30}
\definecolor{techblue}{RGB}{0,90,160}
\definecolor{figbg}{RGB}{245,248,252}
\definecolor{techbg}{RGB}{240,248,240}
\definecolor{toolbg}{RGB}{252,248,235}

% Section formatting
\titleformat{\section}{\large\bfseries\color{primary}\sffamily}{}{0em}{}[\vspace{-4pt}\textcolor{primary!50}{\hrule height 1pt}\vspace{2pt}]
\titleformat{\subsection}{\normalsize\bfseries\color{accent}}{}{0em}{}[\vspace{-2pt}\textcolor{accent!60}{\hrule height 0.5pt}\vspace{2pt}]

% Figure entry box
\newtcolorbox{figbox}[2]{
  title={\textbf{\textsf{#1}}\enspace\textit{\small #2}},
  colback=figbg, colframe=primary!40,
  fonttitle=\small, coltitle=primary,
  arc=3pt, boxrule=0.7pt,
  top=5pt, bottom=5pt, left=7pt, right=7pt,
  before skip=5pt, after skip=5pt,
  breakable
}

% Tool entry box
\newtcolorbox{toolbox}[2]{
  title={\textbf{#1}\enspace\textsf{\small --- #2}},
  colback=toolbg, colframe=techblue!50,
  fonttitle=\small, coltitle=techblue,
  arc=3pt, boxrule=0.7pt,
  top=5pt, bottom=5pt, left=7pt, right=7pt,
  before skip=5pt, after skip=5pt,
  breakable
}

\pagestyle{fancy}
\fancyhf{}
\fancyhead[L]{\small\textcolor{primary}{\textbf{AI NéphroCare}} --- Explication détaillée des figures}
\fancyhead[R]{\small\thepage}
\renewcommand{\headrulewidth}{0.4pt}

\hypersetup{colorlinks=true, linkcolor=primary, pdftitle={AI NéphroCare — Figures et outils}}

%==========================================================
\begin{document}
%==========================================================

\begin{titlepage}
\centering\vspace*{1.5cm}
{\LARGE\bfseries\color{primary}AI NéphroCare\\[0.5cm]}
{\Large\bfseries Explication détaillée de chaque figure,\\
diagramme et outil du rapport\\[0.8cm]}
{\normalsize\color{gray}
Plateforme d'aide à la décision clinique en hémodialyse\\
Cohorte HD-478 $\cdot$ CHU Hassan II de Fès $\cdot$ 2020--2024\\[0.3cm]
SVM Linéaire $\cdot$ AUC = 0{,}826 $\cdot$ 478 patients $\cdot$ 95 décès à 1 an\\[1cm]}
{\small
\begin{tabular}{ll}
\textbf{Modèle ML :} & SVM Linéaire, C=0{,}01, 32 features (24 cliniques + 8 interactions)\\
\textbf{Pipeline :} & KNNImputer(k=3) $\rightarrow$ PowerTransformer(Yeo-Johnson) $\rightarrow$ SVC\\
\textbf{Calibration :} & CalibratedClassifierCV (Platt/sigmoid, cv=5)\\
\textbf{Backend :} & Django REST Framework 4.2, PostgreSQL 16, Redis 7, Celery\\
\textbf{LLM :} & Qwen2.5:14B via Ollama, ChromaDB (RAG), 7 services Docker\\
\end{tabular}}
\vfill
\end{titlepage}

\tableofcontents\newpage

%==========================================================
\section{Outils et technologies --- rôle précis dans le projet}
%==========================================================

Cette section explique \textbf{exactement ce que fait chaque outil} dans le projet
AI NéphroCare, avec les fichiers et paramètres réels.

\begin{toolbox}{Django REST Framework 4.2}{Backend API --- 5 applications}
\textbf{Rôle dans le projet :} Sert les 40+ endpoints REST qui font fonctionner
toute la plateforme. Le backend est organisé en 5 applications Django~:

\begin{itemize}[leftmargin=1.5em, itemsep=2pt]
  \item \texttt{users} --- gestion des comptes, authentification JWT personnalisée.
    \texttt{CustomTokenSerializer} surcharge \texttt{get\_token()} pour injecter
    dans le payload~: \texttt{id}, \texttt{email}, \texttt{nom}, \texttt{prenom},
    \texttt{telephone}, \texttt{role}.
  \item \texttt{patients} --- modèle \texttt{Patient} avec 160+ champs répartis en
    11 sections JSONB~; \texttt{save()} auto-génère \texttt{id\_patient}~;
    \texttt{icc\_charlson} est un \texttt{CharField} (pas Integer)~; pipeline LLM
    de prétraitement complet dans \texttt{views.py} (3800+ lignes).
  \item \texttt{predictions} --- charge \texttt{mortalite\_svm.joblib}, extrait
    les 32 features via \texttt{extract\_features\_for\_patient()}, appelle
    \texttt{pipeline.predict\_proba(input\_df)[:,1][0]}, sauve en JSON filesystem
    via \texttt{\_save\_last\_prediction()}.
  \item \texttt{audit} --- \texttt{AuditLog} avec \texttt{action=TextField}~;
    chaque opération sensible y est enregistrée automatiquement.
  \item \texttt{config} --- \texttt{urls.py} routant 5 préfixes API~:
    \texttt{/api/auth/}, \texttt{/api/patients/}, \texttt{/api/predictions/},
    \texttt{/api/audit/}.
\end{itemize}
\textbf{RBAC :} classes \texttt{IsSuperAdmin} et \texttt{IsAdminOrChefService}
dans \texttt{users/permissions.py}.
\end{toolbox}

\begin{toolbox}{React 18 + Axios}{Frontend SPA --- 6 pages}
\textbf{Rôle dans le projet :} Application monopage (SPA) avec 6 pages principales~:
Login, Dashboard, Patients, Prétraitement, Modèle IA, Monitor.

\begin{itemize}[leftmargin=1.5em, itemsep=2pt]
  \item \texttt{AuthContext.js} --- après login réussi, extrait
    \texttt{\{access, refresh\}} de la réponse, stocke dans localStorage, décode
    le JWT via \texttt{jwtDecode(access)} pour initialiser l'état utilisateur
    (id, email, nom, rôle). Pas de refresh automatique sur 401.
  \item \texttt{axios.js} --- intercepteur REQUEST uniquement~: ajoute
    \texttt{Authorization: Bearer <token>} à chaque requête. Aucun intercepteur
    RESPONSE (pas d'auto-refresh sur expiration).
  \item \texttt{ProtectedRoute} --- filtre la navigation selon le rôle RBAC~;
    les sections non autorisées ne sont jamais rendues.
  \item Chart.js --- rend les 12 graphiques analytiques (6 patients + 6 modèle IA).
\end{itemize}
\end{toolbox}

\begin{toolbox}{PostgreSQL 16}{Base de données --- 9 tables, port 5432}
\textbf{Rôle dans le projet :} Stocke les données persistantes dans la base
\texttt{plateforme\_medicale}.

\begin{itemize}[leftmargin=1.5em, itemsep=2pt]
  \item 9 tables~: \texttt{users\_role}, \texttt{users\_utilisateur},
    \texttt{patients\_patient}, \texttt{patients\_preprocessvalidationrequest},
    \texttt{patients\_dynamiccolumnrequest}, \texttt{patients\_patientformtemplate},
    \texttt{patients\_patientformfield}, \texttt{audit\_auditlog},
    \texttt{audit\_hiddenauditlog}.
  \item \texttt{patients\_patient} contient 160+ colonnes dont 11 colonnes JSONB
    (une par section clinique) + \texttt{extra\_data JSONB} pour les colonnes
    dynamiques importées depuis Excel.
  \item Index GIN sur les colonnes JSONB pour accélérer les recherches dans les
    sections cliniques.
  \item Données persistées via volume Docker \texttt{postgres\_data}.
\end{itemize}
\end{toolbox}

\begin{toolbox}{Redis 7}{Broker Celery --- port 6379}
\textbf{Rôle dans le projet :} Redis est utilisé \textbf{exclusivement} comme
broker de tâches Celery et backend de résultats~:
\begin{itemize}[leftmargin=1.5em, itemsep=2pt]
  \item \texttt{CELERY\_BROKER\_URL = 'redis://redis:6379/0'}
  \item \texttt{CELERY\_RESULT\_BACKEND = 'redis://redis:6379/0'}
\end{itemize}
\textbf{Important :} Redis ne stocke \textit{pas} les sessions de prétraitement.
Celles-ci sont sauvegardées sous forme de fichiers JSON sur le filesystem du
conteneur backend (\texttt{patients/preprocess\_sessions/\{id\}.json}).
\end{toolbox}

\begin{toolbox}{Celery Worker}{Traitement asynchrone du LLM}
\textbf{Rôle dans le projet :} Exécute la tâche \texttt{analyze\_preprocess\_async}
définie dans \texttt{patients/tasks.py} pour éviter les timeouts HTTP pendant les
5--10 minutes d'inférence Qwen2.5:14B.

Paramètres de la tâche~: \texttt{session\_id} (UUID hex), \texttt{file\_path}
(chemin temporaire dans \texttt{/tmp/preprocess/}), \texttt{user\_id},
\texttt{use\_llm} (booléen). Le pipeline exécuté~:
\begin{enumerate}[leftmargin=2em, itemsep=1pt]
  \item Lecture du fichier Excel/CSV via Pandas
  \item Profilage technique (\texttt{\_build\_technical\_profile()})
  \item Chunking sémantique (\texttt{\_build\_preprocess\_chunks()})
  \item RAG ChromaDB (\texttt{\_build\_retrieval\_context()})
  \item Appel LLM Qwen2.5:14B (\texttt{\_call\_ollama\_qwen\_analysis()})
  \item Correction KB médicale (\texttt{\_run\_bio\_value\_correction\_pass()})
  \item Imputation KNN (\texttt{\_run\_model\_imputation()}, k=5 ici)
  \item Sauvegarde session JSON (status=``completed'')
\end{enumerate}
\end{toolbox}

\begin{toolbox}{Ollama + Qwen2.5:14B}{LLM local --- port 11434}
\textbf{Rôle dans le projet :} Modèle de langage 14 milliards de paramètres,
déployé localement sans API externe, pour analyser les fichiers de données
cliniques et proposer des corrections justifiées médicalement.

\begin{itemize}[leftmargin=1.5em, itemsep=2pt]
  \item Endpoint appelé~: \texttt{POST \{OLLAMA\_BASE\_URL\}/api/generate}
    (défaut Docker~: \texttt{http://ollama:11434/api/generate})
  \item Entrée~: prompt structuré contenant le profil technique du fichier
    (nb lignes, colonnes, \% manquants, valeurs aberrantes) + contexte médical
    RAG (normes néphrologiques de ChromaDB) + échantillons de données
  \item Sortie attendue~: JSON avec \texttt{corrections}, \texttt{issues},
    \texttt{pipeline}, \texttt{route}, \texttt{confidence\_contract}
  \item Justification du choix~: déploiement on-premise (confidentialité des
    données patients), pas de coût d'API, contrôle total sur le modèle
\end{itemize}
\end{toolbox}

\begin{toolbox}{ChromaDB}{Base vectorielle RAG --- stockage local}
\textbf{Rôle dans le projet :} Stocke et recherche des embeddings de chunks
cliniques pour fournir un contexte médical pertinent au LLM (RAG).

\begin{itemize}[leftmargin=1.5em, itemsep=2pt]
  \item Mode~: \texttt{PersistentClient(path='/tmp/preprocess/chroma')}
    (client embarqué, pas un serveur séparé)
  \item Modèle d'embedding~: \texttt{nomic-embed-text} via Ollama
  \item Nommage des collections~: hash SHA-1 des métadonnées du fichier
    (lignes, colonnes, \% manquants) pour éviter les doublons
  \item Fallback~: si ChromaDB est indisponible, la fonction \texttt{build\_medical\_rag\_context()} utilise
    \texttt{medical\_kb.json} (base de connaissances locale) pour fournir
    les normes néphrologiques (albumine, calcium, PTH, ferritine, etc.)
  \item Requête~: \texttt{collection.query(query\_embeddings=..., n\_results=max\_chunks,
    include=['documents','metadatas','distances'])}
\end{itemize}
\end{toolbox}

\begin{toolbox}{scikit-learn}{Pipeline ML --- SVM + calibration}
\textbf{Rôle dans le projet :} Implémente le modèle de prédiction de mortalité
à 1 an. Pipeline exact~:

\texttt{Pipeline(steps=[('imputer', KNNImputer(n\_neighbors=3)), ('transformer',
PowerTransformer(method='yeo-johnson')), ('clf', SVC(C=0.01, kernel='linear',
class\_weight='balanced', probability=True, random\_state=42))])}

Enveloppé dans~:
\texttt{CalibratedClassifierCV(base\_pipe, method='sigmoid', cv=5)} = Platt scaling.

\begin{itemize}[leftmargin=1.5em, itemsep=2pt]
  \item \textbf{KNNImputer(k=3)~:} impute les valeurs manquantes des 32 features
    par moyenne des 3 voisins les plus proches dans l'espace des features
    disponibles. Indispensable car 8/32 features ont des valeurs manquantes dans
    HD-478 (ex.~: \texttt{du\_residuelle} manque pour 40\% des patients).
  \item \textbf{PowerTransformer(Yeo-Johnson)~:} normalise la distribution de
    chaque feature vers une gaussienne, quelle que soit son signe (contrairement
    à Box-Cox qui exige des valeurs positives). Indispensable pour SVM qui est
    sensible aux échelles.
  \item \textbf{SVC(C=0{,}01, linear, balanced)~:} marge large (petit C) pour
    éviter le sur-apprentissage sur les 478 patients~; \texttt{class\_weight=balanced}
    compense le déséquilibre 383 survivants / 95 décès (80/20).
  \item \textbf{Calibration Platt~:} transforme les scores de décision SVC en
    probabilités calibrées. Validé par le Brier Score = 0{,}130 (excellent~$<$~0{,}15).
\end{itemize}
Sauvegardé dans \texttt{predictions/models/mortalite\_svm.joblib} (compress=3).
Métadonnées dans \texttt{mortalite\_svm\_features.joblib}~:
feature\_keys (32), threshold (Youden), metrics (AUC, PR-AUC, Brier, F1).
\end{toolbox}

\begin{toolbox}{SimpleJWT}{Authentification --- HS256, 8h/1j}
\textbf{Rôle dans le projet :} Gère l'émission et la vérification des tokens JWT.
Configuration dans \texttt{config/settings.py}~:
\texttt{ACCESS\_TOKEN\_LIFETIME = timedelta(hours=8)},
\texttt{REFRESH\_TOKEN\_LIFETIME = timedelta(days=1)},
\texttt{ALGORITHM = 'HS256'}, \texttt{AUTH\_HEADER\_TYPES = ('Bearer',)}.

Payload custom (ajouté par \texttt{CustomTokenSerializer.get\_token()})~:
\texttt{id}, \texttt{email}, \texttt{nom}, \texttt{prenom}, \texttt{telephone},
\texttt{role}. Plus les claims standard SimpleJWT~: \texttt{user\_id},
\texttt{exp}, \texttt{iat}, \texttt{jti}, \texttt{token\_type}.

Endpoints~: \texttt{POST /api/auth/login/} et
\texttt{POST /api/auth/token/refresh/}.
\end{toolbox}

\begin{toolbox}{Docker Compose}{Infrastructure --- 7 services}
\textbf{Rôle dans le projet :} Orchestre tous les composants en un seul
\texttt{docker compose up --build -d}. Les 7 services~:
\begin{enumerate}[leftmargin=2em, itemsep=1pt]
  \item \textbf{medical\_db} --- PostgreSQL 16, volume \texttt{postgres\_data},
    healthcheck avant démarrage du backend
  \item \textbf{medical\_redis} --- Redis 7, broker Celery
  \item \textbf{medical\_backend} --- Django, dépend de db + redis healthy
  \item \textbf{medical\_frontend} --- React build statique servi par serve/Nginx
  \item \textbf{medical\_ollama} --- Qwen2.5:14B, port 11434, inférence GPU/CPU
  \item \textbf{medical\_celery\_worker} --- même image que backend, commande
    \texttt{celery -A config worker}
  \item \textbf{medical\_audit\_cleanup} --- tâche périodique de nettoyage des
    vieux logs d'audit
\end{enumerate}
\end{toolbox}

\begin{toolbox}{openpyxl}{Export Excel coloré}
\textbf{Rôle dans le projet :} Génère le fichier Excel exporté après prétraitement
LLM avec mise en forme visuelle~:
\begin{itemize}[leftmargin=1.5em, itemsep=2pt]
  \item Feuille \textit{Données}~: cellules modifiées par le LLM en jaune pâle
    (\texttt{\#FFF2CC}) avec police orange gras --- identifiable d'un coup d'œil
  \item Feuille \textit{Corrections}~: tableau récapitulatif de chaque modification
    (valeur originale, valeur corrigée, type de correction, justification médicale)
  \item Généré par \texttt{PatientPreprocessExportView} sur
    \texttt{GET /api/patients/preprocess/\{id\}/export/}
\end{itemize}
\end{toolbox}

\begin{toolbox}{Pandas}{Profiling et manipulation des données}
\textbf{Rôle dans le projet :} Lit et profile les fichiers Excel/CSV importés~:
\texttt{pd.read\_excel()}, \texttt{pd.read\_csv()} avec détection automatique
du séparateur. Normalise les noms de colonnes (lowercase, strip). Calcule le
profil technique~: dtype par colonne, \% valeurs manquantes, nb de doublons,
valeurs aberrantes (IQR). Construit les \texttt{DataFrame} d'entrée pour le
pipeline SVM (\texttt{pd.DataFrame([row], columns=feature\_keys)}).
\end{toolbox}

\begin{toolbox}{SHAP + XGBoost}{Sélection des features --- analyse post-hoc}
\textbf{Rôle dans le projet :} Utilisé en dehors de la plateforme pour justifier
la sélection des 24 variables parmi 77 candidates sur HD-478. XGBoost sert de
modèle de référence pour calculer les valeurs SHAP (Shapley Additive
exPlanations)~: chaque variable reçoit un score d'importance mesurant sa
contribution marginale moyenne à la prédiction. Les 24 variables avec les
valeurs SHAP les plus élevées (top-24 sur 77) constituent le sous-ensemble
sélectionné pour le SVM, auquel s'ajoutent 8 termes d'interaction biologiquement
justifiés --- soit \textbf{32 features} au total.
\end{toolbox}

\begin{toolbox}{Chart.js}{Visualisation --- 12 graphiques interactifs}
\textbf{Rôle dans le projet :} Rend les 12 graphiques dynamiques de la plateforme~:
\textbf{Module Patients (6)}~: tendance mensuelle des inclusions, camembert sexe,
histogramme âge, barres complications, barres étiologies, barres comorbidités.
\textbf{Module IA (6)}~: camembert zones de risque, barres issues prédites,
courbe Kaplan-Meier (survie cumulée), barres facteurs contributifs, radar profil
de risque, boxplots biologiques par zone. Tous calculés à la volée sur la cohorte
active et interactifs (hover, zoom).
\end{toolbox}

\newpage
%==========================================================
\section{Chapitre 1 --- Introduction}
%==========================================================

\begin{figbox}{Figure 1.1}{Logo FMPDF}
Logotype officiel de la Faculté de Médecine, de Pharmacie et de Dentisterie de
Fès (FMPDF), établissement de tutelle académique du projet. Sa présence dans le
rapport situe institutionnellement le travail dans le système universitaire
marocain (Université Sidi Mohamed Ben Abdellah, Fès).
\end{figbox}

\begin{figbox}{Figure 1.2}{Logo Centre HAIM}
Logotype du \textit{Health Artificial Intelligence \& Modeling Center} (HAIM),
structure de recherche hébergeant le développement de la plateforme. Le HAIM
opère en partenariat avec le service de Néphrologie du CHU Hassan II de Fès,
qui a fourni les données de la cohorte HD-478 (478 patients hémodialysés,
2020--2024, 95 décès à 1 an, soit un taux de mortalité annuel de 19{,}9\%).
\end{figbox}

\begin{figbox}{Figure 1.3}{Diagramme Kanban Trello --- 8 colonnes, 42 tâches}
Le diagramme représente le tableau Trello utilisé pour piloter le projet en
mode Agile sur 9 sprints de 2 semaines (octobre 2024 -- juin 2025). Les 8
colonnes correspondent au cycle de vie complet de chaque tâche~:
\textbf{Backlog} (idées non planifiées), \textbf{In Progress} (développement
actif), \textbf{Review} (relecture code/résultats), \textbf{To Validate} (en
attente de validation par le superviseur technique M. Zakaria Alami),
\textbf{Done/Validated} (livrés et intégrés). Les cartes en orange (In Progress)
incluent les tâches lourdes~: pipeline SVM, intégration ChromaDB, tâches Celery.
Les cartes vertes (Done) incluent les 7 services Docker, le schéma PostgreSQL
(9 tables), le SPA React (6 pages), le pipeline LLM complet, et le SVM
(AUC~=~0{,}826).
\end{figbox}

\begin{figbox}{Figure 1.4}{Stadification IRC selon KDIGO 2012}
Grille officielle KDIGO 2012 croisant les 5 stades de DFGe (G1~:~$\geq$90,
G2~: 60--89, G3a~: 45--59, G3b~: 30--44, G4~: 15--29, G5~:~$<$15 mL/min/1{,}73m²)
avec les 3 catégories d'albuminurie (A1~:~$<$30, A2~: 30--300, A3~:~$>$300 mg/g).
Le code couleur encode le risque composite (vert~=~faible, jaune~=~modéré,
orange~=~élevé, rouge~=~très élevé). Les patients HD-478 sont tous en stade
\textbf{G5D} (dialyse, DFGe~$<$15), zone rouge foncé, justifiant la concentration
du modèle sur la mortalité à 1 an comme critère primaire.
\end{figbox}

%==========================================================
\section{Chapitre 2 --- État de l'art}
%==========================================================

\begin{figbox}{Tableau 2.1}{Comparaison des études ML de prédiction de mortalité HD à 1 an}
Ce tableau compare \textbf{7 études internationales publiées entre 2016 et 2023},
toutes avec le même endpoint (mortalité toutes causes à 1 an en hémodialyse),
ce qui rend les AUC-ROC directement comparables. Colonnes~: auteurs/année,
cohorte ($n$), pays, algorithmes testés, AUC meilleur, top variables.

Valeurs clés issues des études~:
\begin{itemize}[leftmargin=1.5em, itemsep=2pt]
  \item \textbf{Sheng 2020} (PMID~33118948, JMIR Medical Informatics)~:
    $n$~=~1825, XGBoost, AUC~=~0{,}843, 42 features candidates, top-15 incluent
    accès vasculaire, albumine, calcium, PTH, ferritine, cardiopathie, diabète
    néphropathique. \textit{Référence principale} pour la sélection des 24 variables.
  \item \textbf{Fotheringham 2022}~: $n$~=~56\,000 (USRDS), AUC~=~0{,}79--0{,}84.
  \item \textbf{Notre modèle SVM HD-478}~: $n$~=~478, AUC~=~0{,}826 --- dans la
    médiane haute de l'état de l'art malgré la taille de cohorte 4$\times$ plus
    petite, validant la rigueur du feature engineering.
\end{itemize}
\end{figbox}

%==========================================================
\section{Chapitre 3 --- Vue d'ensemble}
%==========================================================

\begin{figbox}{Figure 3.1}{Diagramme des cas d'utilisation UML}
Diagramme UML avec \textbf{4 acteurs} et \textbf{12 cas d'utilisation}, montrant
qui peut faire quoi dans la plateforme.

Les 4 acteurs et leurs droits spécifiques~:
\begin{itemize}[leftmargin=1.5em, itemsep=2pt]
  \item \textbf{Super Admin}~: tous droits~; créer/supprimer comptes de tous niveaux~;
    voir tous les AuditLogs~; approuver/rejeter imports~; prédire mortalité.
  \item \textbf{Chef de Service}~: comme Super Admin sauf~: ne peut pas gérer
    les comptes Super Admin~; peut approuver les imports.
  \item \textbf{Professeur}~: consulter patients, importer fichiers (en attente
    d'approbation), prédire mortalité, voir ses propres activités.
  \item \textbf{Résident}~: comme Professeur~; accès restreint aux statistiques
    avancées.
\end{itemize}
Les généralisations (flèches d'héritage) indiquent que chaque acteur hérite des
droits du niveau inférieur. Ce diagramme sert de base au RBAC implémenté dans
\texttt{users/permissions.py}.
\end{figbox}

\begin{figbox}{Figure 3.2 et Tableau 3.1}{Architecture en cinq couches}
Architecture logicielle à 5 couches superposées~:
\begin{enumerate}[leftmargin=2em, itemsep=2pt]
  \item \textbf{Présentation} --- React 18 SPA, Material-UI, Chart.js, Axios.
    Communication~: HTTP/HTTPS sur port 3000.
  \item \textbf{Application} --- Django REST Framework 4.2, port 8000.
    Reçoit toutes les requêtes API, applique RBAC, coordonne les composants.
  \item \textbf{IA/ML} --- SVM joblib (prédiction synchrone), Qwen2.5:14B via
    Ollama (prétraitement asynchrone via Celery), ChromaDB (RAG embeddings).
  \item \textbf{Données} --- PostgreSQL 16 (port 5432, données permanentes),
    Redis 7 (port 6379, broker Celery), Filesystem JSON (sessions prétraitement).
  \item \textbf{Infrastructure} --- Docker Compose 7 services,
    variables d'environnement \texttt{.env}, healthchecks entre services.
\end{enumerate}
Le tableau 3.1 précise les protocoles entre couches~: REST JSON (1$eftrightarrow$2),
ORM Django/SQL (2$eftrightarrow$4), urllib HTTP (2$eftrightarrow$3 Ollama), joblib (2$eftrightarrow$3 SVM),
TCP socket (2$eftrightarrow$4 Redis).
\end{figbox}

\begin{figbox}{Figure 3.3}{Stack technologique complet}
Vue d'ensemble de \textbf{toutes les technologies} regroupées par domaine.
Ce diagramme répond à la question ``quels outils sont utilisés ensemble~?''~:
Django (backend) + PostgreSQL (données) + Redis (queue) + Celery (async) +
Ollama (LLM) + ChromaDB (RAG) + React (UI) + Chart.js (graphiques) +
scikit-learn (ML) + Docker (déploiement) --- 14 technologies coordonnées dans
un seul \texttt{docker compose up}.
\end{figbox}

%==========================================================
\section{Chapitre 4 --- Pipeline de données}
%==========================================================

\subsection{Curation manuelle HD-478}

\begin{figbox}{Tableau 4.1}{Dataset HD-478 avant vs.\ après curation}
Comparaison quantitative du fichier Excel brut hospitalier vs.\ dataset
ML-ready après curation manuelle (9 catégories de corrections)~:

\begin{center}
\begin{tabular}{lcc}
\toprule
\textbf{Indicateur} & \textbf{Original} & \textbf{Corrigé} \\
\midrule
Patients & 478 & 478 \\
Variables & 160+ & 160+ \\
Valeurs manquantes & $\sim$34\% & $\sim$18\% \\
Valeurs aberrantes biologiques & 287 & 0 \\
Numéros sériels Excel (dates) & 312 & 0 (dates ISO) \\
DFGe incohérents & 109 & 0 (recalculés MDRD) \\
Marqueurs X/x & 265 & 0 (0 ou 1) \\
\bottomrule
\end{tabular}
\end{center}
La curation a réduit les valeurs manquantes de 34\% à 18\% et éliminé tous les
artéfacts de saisie, permettant l'entraînement d'un modèle ML fiable.
\end{figbox}

\begin{figbox}{Tableau 4.2}{Conversion dates Excel (312 valeurs)}
Le format Excel stocke les dates comme des entiers (numéros sériels depuis le
01/01/1900). Ex.~: le nombre 44927 correspond au 01/01/2023. Sans conversion,
Pandas lit ces cellules comme des entiers, rendant impossible tout calcul de
survie (délai entre inclusion et décès). La formule appliquée~:
$\text{date} = \text{date\_epoch} + \text{serial} \times 86400\,\text{s}$.
Colonnes corrigées~: \texttt{date\_debut\_dialyse}, \texttt{date\_inclusion},
\texttt{date\_deces} --- essentielles pour calculer
\texttt{delai\_jusquau\_deces\_jours} utilisé par la variable cible du modèle.
\end{figbox}

\begin{figbox}{Tableau 4.3}{Valeurs biologiques aberrantes (287 corrections)}
287 valeurs biologiques aberrantes détectées par règles cliniques spécifiques~:
albumine $>$60 g/L (impossible physiologiquement, max normal 55 g/L)~;
créatinine $>$2000~$\mu$mol/L~; sodium $<$100 ou $>$175~mEq/L~; calcium $<$1 ou $>$4
mmol/L~; PTH $>$5000~pg/mL~; ferritine $>$10\,000~$\mu$g/L. Chaque valeur aberrante
est remplacée par NaN puis imputée par KNNImputer dans le pipeline ML,
ou corrigée manuellement si une explication clinique existe (ex.~: unité erronée~:
albumine en mg/dL au lieu de g/L $\rightarrow$ divisé par 10).
\end{figbox}

\begin{figbox}{Tableau 4.4}{Recalcul DFGe MDRD (109 valeurs)}
Le DFGe (débit de filtration glomérulaire estimé) de 109 patients était absent
ou incohérent. Recalcul par la formule MDRD-4~:
$\text{DFGe} = 175 \times [\text{créatinine}_{(\mu\text{mol/L})}/88{,}4]^{-1{,}154} \times
\text{âge}^{-0{,}203} \times [0{,}742\text{ si femme}]$.
Le DFGe est une variable-clé pour confirmer le stade G5D (DFGe~$<$15) et
n'est pas utilisée directement dans le modèle SVM (évite la fuite de données),
mais sert à la validation clinique des dossiers.
\end{figbox}

\begin{figbox}{Tableau 4.5}{Nettoyage marqueurs X/x --- 265 corrections}
Dans les colonnes d'éducation (\texttt{niveau\_education}) et de statut
transplantation, les soignants avaient coché ``X'' ou ``x'' au lieu d'écrire
0 ou 1. Python/Pandas interprétait ces cellules comme des \texttt{object}
(chaîne de caractères), empêchant tout calcul statistique ou ML. Chaque ``X''/``x''
a été remplacé par la valeur numérique correspondant à la sémantique de la
colonne (0~=~absent, 1~=~présent). 265 corrections réparties sur 8 colonnes.
\end{figbox}

\begin{figbox}{Tableau 4.6}{Imputation NaN $\rightarrow$ 0 pour variables transplantation}
Pour 3 variables binaires (\texttt{liste\_attente\_transplantation},
\texttt{information\_transplantation\_donnee}, \texttt{transplantation\_effectuee})~:
les valeurs manquantes signifient cliniquement ``non'' (absence de mention dans
le dossier = événement non survenu). Imputation $\text{NaN} \rightarrow 0$
appliquée à 187 cellules. Cette règle logique est préférable à l'imputation
statistique qui introduirait un biais~: imputer 0{,}3 pour
\texttt{liste\_attente\_transplantation} n'a pas de sens clinique.
\end{figbox}

\begin{figbox}{Tableau 4.7}{Distribution des types de variables (dataset corrigé)}{%
Après curation, les 160+ variables se répartissent~: 68\% binaires (0/1,
ex.~: comorbidités, complications), 24\% numériques continues (biologie,
durées), 8\% catégorielles multi-classes (étiologie, modalité). Cette
distribution justifie l'usage de \texttt{class\_weight=balanced} dans le SVM
(les variables binaires ne dominent pas, mais la cible est déséquilibrée
80/20) et le choix de \texttt{PowerTransformer} (normalise les distributions
asymétriques des variables biologiques continues).}
\end{figbox}

\subsection{Schéma de base de données}

\begin{figbox}{Figure 4.1}{Diagramme ER --- schéma complet PostgreSQL}
Diagramme Entité-Relations des \textbf{9 tables} de \texttt{plateforme\_medicale}.
Relations clés~:
\begin{itemize}[leftmargin=1.5em, itemsep=2pt]
  \item \texttt{users\_utilisateur} $\rightarrow$ \texttt{users\_role} (FK,
    chaque utilisateur a un rôle parmi 4~: super\_admin, chef\_service,
    professeur, résident)
  \item \texttt{patients\_patient.utilisateur\_saisie} est un \textbf{CharField}
    (pas une FK) --- stocke le label de l'utilisateur au moment de la saisie
    pour préserver l'historique si le compte est supprimé
  \item \texttt{patients\_preprocessvalidationrequest} $\rightarrow$
    \texttt{users\_utilisateur} (submitted\_by FK + reviewed\_by FK, nullable)
  \item \texttt{audit\_auditlog.action} est un \textbf{TextField} (pas CharField)
    pour accueillir des descriptions longues sans troncature
  \item \texttt{audit\_hiddenauditlog} $\rightarrow$ \texttt{audit\_auditlog}
    (pour masquer certaines entrées d'audit sans les supprimer)
\end{itemize}
\end{figbox}

\begin{figbox}{Tableaux 4.8 à 4.12}{Schémas détaillés des tables PostgreSQL}
\textbf{Tableau 4.8 --- \texttt{users\_role} et \texttt{users\_utilisateur}~:}
\texttt{users\_role}~: id, nom (CharField max=50, UNIQUE, 4 choix).
\texttt{users\_utilisateur}~: hérite AbstractBaseUser~; email UNIQUE~; nom, prénom,
téléphone CharField~; personal\_notes TextField~; role FK~; is\_active BooleanField~;
date\_creation auto\_now\_add.

\textbf{Tableau 4.9 --- \texttt{patients\_patient}~:}
160+ champs~; icc\_charlson~=~\texttt{CharField(max\_length=10)}~;
\texttt{save()} override auto-génère \texttt{id\_patient} (format PAT-XXXXXX)~;
11 colonnes JSONB pour les sections cliniques~; \texttt{extra\_data JSONB} pour
les colonnes dynamiques.

\textbf{Tableau 4.10 --- \texttt{patients\_preprocessvalidationrequest}~:}
session\_id CharField(64), status (pending/approved/rejected), submitted\_by FK,
reviewed\_by FK (nullable), quality\_score FloatField, source\_file\_name,
submitted\_at, reviewed\_at, comment, rows\_count, columns\_count.

\textbf{Tableau 4.11 --- \texttt{audit\_auditlog}~:}
action TextField, entite CharField, entite\_id IntegerField, details TextField,
adresse\_ip GenericIPAddressField, date DateTimeField, utilisateur FK.

\textbf{Tableau 4.12 --- Sections cliniques du modèle Patient~:}
11 sections couvrent le parcours de soin néphologique complet~:
Identité/démographie, Antécédents IRC, Comorbidités, Biologie basale,
Dialyse, Complications, Médicaments, Transplantation, Suivi, Outcomes,
Données dynamiques (extra\_data JSONB).
\end{figbox}

\subsection{Pipeline LLM de prétraitement}

\begin{figbox}{Figure 4.2}{Pipeline LLM de prétraitement automatisé}
Ce diagramme est le plus important du chapitre 4~: il montre comment un fichier
Excel brut devient un dataset clinique valide sans saisie manuelle.

\textbf{11 étapes avec acteurs réels~:}
\begin{enumerate}[leftmargin=2em, itemsep=2pt]
  \item Clinicien uploade le fichier via React~: \texttt{POST /api/patients/preprocess/analyze/}
  \item Django génère un \texttt{session\_id} UUID hex, sauve le fichier dans
    \texttt{/tmp/preprocess/}, crée la session JSON (status=``pending'')
  \item Django retourne \texttt{\{preprocess\_id: "<uuid>"\}} au frontend
  \item Django dispatche \texttt{analyze\_preprocess\_async.delay()} vers Celery
  \item \textbf{Celery Worker (asynchrone)~:} Pandas profile le fichier
    (nb lignes, colonnes, \% manquants, dtypes, outliers)
  \item Celery appelle \texttt{\_build\_retrieval\_context()}~: ChromaDB
    (PersistentClient) génère embeddings via \texttt{nomic-embed-text} et
    récupère les normes néphrologiques pertinentes
  \item Celery envoie \texttt{POST http://ollama:11434/api/generate} avec
    prompt structuré (profil technique + contexte RAG + échantillons de données)
  \item Qwen2.5:14B retourne un JSON de corrections avec justifications
  \item Celery applique 3 passes de correction~: plan LLM, KB biologique
    (normes locales \texttt{medical\_kb.json}), imputation KNN (k=5)
  \item Session sauvegardée en JSON filesystem (status=``completed'')
  \item Frontend polle \texttt{GET /api/patients/preprocess/\{id\}/rows/}~;
    clinicien révise~; Chef approuve via
    \texttt{POST /api/patients/preprocess/\{id\}/integrate/}~: INSERT PostgreSQL + AuditLog
\end{enumerate}
\end{figbox}

\subsection{Feature Engineering}

\begin{figbox}{Tableau 4.14}{24 variables cliniques brutes --- domaines et littérature}
Ce tableau relie chaque variable au modèle SVM à une évidence publiée vérifiée.
Groupes~:
\begin{itemize}[leftmargin=1.5em, itemsep=2pt]
  \item \textbf{SHAP \#1--4 (top confirmés)~:} \texttt{fistule\_arterioveineuse\_creee}
    (accès vasculaire, Sheng 2020 top-15), \texttt{couverture\_medicale} (accès aux soins,
    Liyanage 2015 + USRDS 2022), \texttt{calcium\_basale} (Ca dysrégulé prédit
    la mortalité, DOPPS~: Tentori 2008 + Sheng 2020), \texttt{albumine\_basale}
    (chaque $-$1 g/dL~=~risque accru, NDT 2005 + Sheng 2020).
  \item \textbf{SHAP top-24~:} 9 autres variables confirmées SHAP~:
    événements cardiovasculaires, cathéter tunnellisé, diurèse résiduelle,
    hospitalisations, PTH, ferritine, séances/semaine, étiologie MRC, année inclusion.
  \item \textbf{Littérature (11 variables)~:} HTA, cardiopathie, diabète, dyspnée,
    œdèmes, douleur abdominale, crise convulsive, maladie héréditaire rénale,
    hémodialyse, liste attente, information transplantation.
\end{itemize}
\textbf{Toutes les 9 références bibliographiques ont été vérifiées sur PubMed.}
\end{figbox}

\begin{figbox}{Tableau 4.15}{8 termes d'interaction biologiquement justifiés}
Le vecteur de features est étendu de 24 à \textbf{32 dimensions} par 8 termes
d'interaction~:
\begin{center}
\begin{tabular}{ll}
\toprule
\textbf{Interaction} & \textbf{Justification clinique} \\
\midrule
\texttt{age\_x\_charlson} = âge $\times$ ICC & Risque multiplicatif âge/comorbidités\\
\texttt{dyspnee\_x\_oedeme} = dyspnée $\times$ œdèmes & Surcharge hydrique synergique\\
\texttt{sodium\_lt130} = 1 si Na~$<$130 mEq/L & Hyponatrémie sévère binaire\\
\texttt{charlson\_x\_hosp} = ICC $\times$ hospitalisations & Fardeau comorbidités+réadmissions\\
\texttt{hypert\_x\_cardio} = HTA $\times$ cardiopathie & Risque cardiovasculaire cumulé\\
\texttt{albumine\_x\_hosp} = albumine $\times$ hospitalisations & Dénutrition + réadmissions\\
\texttt{albumine\_lt35} = 1 si albumine~$<$35 g/L & Seuil clinique dénutrition\\
\texttt{age\_x\_cardio} = âge $\times$ cardiopathie & Fragilité cardiaque liée à l'âge\\
\bottomrule
\end{tabular}
\end{center}
\end{figbox}

%==========================================================
\section{Chapitre 5 --- Modélisation ML}
%==========================================================

\subsection{Comparaison des modèles}

\begin{figbox}{Tableau 5.1}{Configuration hyperparamètres finaux}
Hyperparamètres retenus après \texttt{RandomizedSearchCV(n\_iter=100, cv=5)}
stratifié pour chacun des 9 algorithmes évalués sur HD-478~:
\begin{center}
\small
\begin{tabular}{ll}
\toprule
\textbf{Modèle} & \textbf{Hyperparamètres clés} \\
\midrule
SVM Linéaire & C=0{,}01, kernel=linear, class\_weight=balanced \\
Régression Log. & C=0{,}1, penalty=l2, solver=lbfgs \\
Random Forest & n\_estimators=200, max\_depth=8, min\_samples\_leaf=3 \\
Extra Trees & n\_estimators=300, max\_features=sqrt \\
Gradient Boosting & n\_estimators=150, learning\_rate=0{,}05, max\_depth=4 \\
AdaBoost & n\_estimators=100, learning\_rate=0{,}1 \\
XGBoost & max\_depth=4, learning\_rate=0{,}05, subsample=0{,}8 \\
LightGBM & num\_leaves=31, learning\_rate=0{,}05 \\
CatBoost & depth=4, learning\_rate=0{,}05, iterations=500 \\
\bottomrule
\end{tabular}
\end{center}
\end{figbox}

\begin{figbox}{Tableau 5.2}{Comparaison des 9 modèles --- résultats HD-478 (10-fold CV)}
Résultats complets de la comparaison~:
\begin{center}
\small
\begin{tabular}{lcccc}
\toprule
\textbf{Modèle} & \textbf{AUC-ROC} & \textbf{PR-AUC} & \textbf{F1} & \textbf{Brier} \\
\midrule
\rowcolor{green!10}\textbf{SVM Linéaire} & \textbf{0{,}826±0{,}042} & \textbf{0{,}573} & \textbf{0{,}512} & \textbf{0{,}130} \\
LDA & 0{,}823±0{,}044 & 0{,}561 & 0{,}498 & 0{,}134 \\
Régression Log. & 0{,}821±0{,}047 & 0{,}558 & 0{,}491 & 0{,}132 \\
Random Forest & 0{,}810±0{,}055 & 0{,}531 & 0{,}487 & 0{,}148 \\
Gradient Boosting & 0{,}805±0{,}061 & 0{,}519 & 0{,}479 & 0{,}152 \\
Extra Trees & 0{,}801±0{,}058 & 0{,}515 & 0{,}471 & 0{,}155 \\
XGBoost & 0{,}809±0{,}053 & 0{,}528 & 0{,}483 & 0{,}149 \\
AdaBoost & 0{,}791±0{,}064 & 0{,}501 & 0{,}462 & 0{,}162 \\
CatBoost & 0{,}803±0{,}057 & 0{,}512 & 0{,}475 & 0{,}153 \\
\bottomrule
\end{tabular}
\end{center}
Le SVM domine sur \textbf{toutes les métriques simultanément}.
\end{figbox}

\begin{figbox}{Figures 5.1--5.2}{Courbes ROC --- 9 modèles + zoom Top 3}
\textbf{Figure 5.1 --- Toutes courbes ROC superposées~:} axe X = taux faux
positifs (1-Spécificité), axe Y = Sensibilité. Diagonale = classifieur aléatoire
(AUC=0{,}5). Le groupe de tête (SVM, LDA, Régression Log.) se détache clairement
des modèles ensemblistes qui présentent une courbure plus irrégulière
(sur-ajustement sur les petits folds).

\textbf{Figure 5.2 --- Zoom Top 3~:} SVM (AUC=0{,}826, courbe la plus haute
dans la région Sensibilité~$>$~0{,}8), LDA (0{,}823), Régression Log. (0{,}821).
La différence est visuellement faible~: le choix final du SVM est confirmé par
les autres métriques (Brier Score, écart train/test, interprétabilité).
\end{figbox}

\begin{figbox}{Figure 5.3}{Barres AUC avec écart-type}
Barres horizontales ordonnées~: AUC moyen $\pm$ 1 écart-type sur les 10 folds.
\textbf{SVM~: $0{,}826 \pm 0{,}042$} --- la barre d'erreur la plus courte parmi
les 3 premiers modèles, confirmant sa stabilité. AdaBoost a le plus grand écart
($\pm 0{,}064$), indicateur d'instabilité aux variations de partition.
\end{figbox}

\begin{figbox}{Figure 5.4}{Écart train-test AUC --- détection sur-apprentissage}
Pour chaque modèle~: AUC\_train (mesuré sur l'ensemble d'entraînement de chaque
fold) vs.\ AUC\_test. Code couleur~: vert si écart~$<$~0{,}05 (excellent),
orange si 0{,}05--0{,}10 (acceptable), rouge si~$>$~0{,}10 (sur-ajustement).

\textbf{SVM~: écart = 0{,}047} (seul modèle dans la zone verte)~; Random Forest~:
0{,}12 (rouge). Cette figure est décisive~: un modèle clinique doit généraliser
sur des patients non vus. Un gap élevé (forêts aléatoires) signifie que le
modèle a mémorisé l'ensemble d'entraînement sans apprendre les patterns généraux.
\end{figbox}

\begin{figbox}{Figure 5.5}{Radar multi-métriques --- 9 modèles}
Graphique radar à \textbf{5 axes}~:
\begin{itemize}[leftmargin=1.5em, itemsep=1pt]
  \item \textbf{AUC-ROC}~: discrimination globale
  \item \textbf{PR-AUC}~: performance sur la classe minoritaire (décès, 20\%)
  \item \textbf{F1-score}~: compromis précision/rappel au seuil Youden
  \item \textbf{Brier Score inversé} ($1-\text{Brier}$)~: qualité de calibration
  \item \textbf{Recall}~: fraction de décès détectés (métrique critique en clinique)
\end{itemize}
Le polygone SVM est le plus \textbf{équilibré} et le plus \textbf{étendu}.
Les forêts aléatoires ont un bon Recall mais un Brier Score élevé (probabilités
mal calibrées). Un bon Recall sans bonne calibration est inutile cliniquement~:
le médecin a besoin de probabilités fiables, pas seulement de classifications.
\end{figbox}

\begin{figbox}{Figure 5.6}{Brier Score par modèle}
Le Brier Score mesure l'erreur quadratique moyenne entre la probabilité prédite
$\hat{p}$ et la réalité $y \in \{0,1\}$~:
$\text{BS} = \frac{1}{n}\sum_{i=1}^{n}(\hat{p}_i - y_i)^2$.
Score parfait~= 0, aléatoire~= 0{,}25 (prévalence 50\%).

\textbf{Seuil excellent~: BS~$<$~0{,}15} (ligne pointillée sur le graphe).
SVM~= 0{,}130 et Régression Log.~= 0{,}132 passent ce seuil, les autres non.
Cliniquement~: si le SVM prédit 65\% de risque, environ 65\% des patients
dans ce décile mourront réellement dans l'année --- les probabilités sont
utilisables pour prioriser les soins.
\end{figbox}

\begin{figbox}{Figure 5.7}{Diagramme de fiabilité (calibration)}
\textbf{Reliability plot} pour les 3 meilleurs modèles. Les probabilités prédites
sont groupées en déciles (axe X~: 0--0{,}1, 0{,}1--0{,}2, ..., 0{,}9--1{,}0).
Pour chaque décile, la fréquence de décès observée est calculée (axe Y).
La diagonale parfaite signifie~: ``si je prédis 30\% de risque, 30\% de ces
patients décèdent effectivement''.

\textbf{SVM + Platt Scaling~:} courbe la plus proche de la diagonale~;
l'écart moyen (\textit{calibration error}) est minimal. Régression Log.~: bien
calibrée aussi mais légèrement optimiste à fort risque. Random Forest~: nettement
sous-calibré aux probabilités élevées (prédit 80\% mais seulement 60\% de décès
observés) --- dangereux cliniquement car surestime la confiance.
\end{figbox}

\begin{figbox}{Figures 5.8 à 5.10}{Matrices de confusion --- SVM, LDA, Régression Log.}
\textbf{Figure 5.8 --- SVM Linéaire} (seuil Youden = 0{,}238)~:
La matrice de confusion présente 4 quadrants~: VP (vrais positifs, décès
correctement détectés), FN (faux négatifs, décès manqués), FP (fausses
alarmes), VN (vrais négatifs, survivants corrects). La disposition est~:
$95$ décès réels (VP + FN) en ligne haute, $383$ survivants réels (FP + VN)
en ligne basse.

Sensibilité (vrais positifs / tous positifs réels)~: fraction des décès correctement
détectés~; Spécificité~: fraction des survivants correctement classés.
La pondération \texttt{class\_weight=balanced} maximise la Sensibilité sur la
classe minoritaire (décès, 20\%). Le seuil de Youden (0{,}238, inférieur au seuil
par défaut de 0{,}5) reflète la nécessité de détecter plus de décès au prix
de quelques fausses alarmes.

\textbf{Figure 5.9 --- LDA} (seuil 0{,}220)~: performances proches du SVM mais
Brier Score légèrement supérieur~; non sélectionné en raison de son hypothèse
de normalité multivariée non vérifiable sur des données réelles mixtes (binaires
+ continues).

\textbf{Figure 5.10 --- Régression Log.} (seuil 0{,}502)~: seuil plus élevé car
distribution de probabilités mieux centrée sur 0{,}5~; bon Brier Score mais
moins stable sur les folds (écart-type AUC = 0{,}047 vs 0{,}042 SVM).
\end{figbox}

\begin{figbox}{Figure 5.11}{Courbes Précision-Rappel}
Les courbes PR sont la \textbf{métrique principale en présence de déséquilibre
de classes} (80\% survivants / 20\% décès). L'axe X est le Rappel (fraction de
décès détectés), l'axe Y la Précision (fraction de vraies alarmes parmi les
alertes émises). La ligne de base~= prévalence~= 0{,}20 (classifieur aléatoire).

PR-AUC SVM~= 0{,}573 (vs.\ 0{,}20 pour le hasard)~: le modèle a une capacité
de détection effective 2{,}9$\times$ supérieure au hasard. Cette figure est
présentée aux cliniciens pour expliquer le compromis~: augmenter le Rappel
(détecter plus de décès) diminue la Précision (plus de fausses alarmes).
\end{figbox}

\begin{figbox}{Figure 5.12}{Boxplot AUC sur 10 folds}
Boîtes à moustaches~: médiane (trait central), Q1-Q3 (boîte), min-max (moustaches).
Ce graphe répond à~: ``est-ce que le modèle est stable ou chanceux~?''

SVM~: boîte étroite centrée sur 0{,}826 --- aucun fold avec AUC~$<$~0{,}78~;
AdaBoost~: boîte large, un fold avec AUC~$<$~0{,}70 (inacceptable en clinique).
La stabilité inter-folds est un critère de confiance clinique~: un modèle qui
varie de 0{,}73 à 0{,}89 selon le fold est imprévisible en production.
\end{figbox}

\subsection{Zones de risque}

\begin{figbox}{Tableau 5.3}{Classification par zones de risque --- seuils HD-478}
Les seuils $T_1$ et $T_2$ sont dérivés par \textbf{bootstrap ROC/Youden}
($n$~=~1000 itérations) sur HD-478 et stockés dans
\texttt{predictions/models/seuils\_classification.joblib}. Chargés
dynamiquement par \texttt{\_resolve\_thresholds(metadata)} avec priorité~:
(1)~\texttt{seuils\_classification.joblib}, (2)~\texttt{calibration\_thresholds},
(3)~\texttt{gmm\_thresholds}, (4)~défaut fixe 0{,}10/0{,}29.

Taux de mortalité observés dans la cohorte HD-478 par zone~:
\begin{center}
\begin{tabular}{lccc}
\toprule
\textbf{Zone} & \textbf{Seuil} & \textbf{Mortalité observée} & \textbf{Patients} \\
\midrule
Faible & $\hat{p} < T_1$ & 3{,}8\% & $\sim$240 \\
Modérée & $T_1 \leq \hat{p} < T_2$ & 15{,}6\% & $\sim$155 \\
Élevée & $\hat{p} \geq T_2$ & 46{,}5\% & $\sim$83 \\
\bottomrule
\end{tabular}
\end{center}
La séparation est significative (log-rank $p<0{,}001$)~: les 3 zones captent
3 populations de survie biologiquement distinctes.
\end{figbox}

\begin{figbox}{Figure 5.13}{Workflow complet de prédiction (diagramme de flux)}
14 étapes représentées avec les acteurs réels du code~:
\begin{enumerate}[leftmargin=2em, itemsep=1pt]
  \item Clic ``Prédire'' sur l'interface React
  \item \texttt{GET /api/predictions/patient/\{id\}/mortalite/}
  \item \texttt{Patient.objects.get(pk=patient\_id)} $	o$ PostgreSQL
  \item Objet ORM retourné
  \item \texttt{svm\_extract\_features(patient)} $	o$ dict de 32 features
  \item \texttt{pd.DataFrame([row], columns=feature\_keys)}
  \item \texttt{pipeline.predict\_proba(input\_df)[:,1][0]} $	o$ $\hat{p}$ calibrée Platt (1 seul appel)
  \item \texttt{\_resolve\_thresholds(metadata)} $	o$ $T_1$, $T_2$ chargés depuis joblib
  \item Classification~: Faible / Modéré / Élevé
  \item \texttt{\_build\_svm\_factors(pipeline, feature\_keys)} $	o$ top-5 coefficients
  \item \texttt{\_save\_last\_prediction(result)} $	o$ JSON filesystem (pas PostgreSQL)
  \item Réponse~: \texttt{\{niveau\_risque, score\_risque, probabilite\_deces,
    factors, recommendation, mort\_rates\}}
  \item Rendu jauge + facteurs dans React
\end{enumerate}
\end{figbox}

%==========================================================
\section{Chapitre 6 --- Architecture logicielle}
%==========================================================

\begin{figbox}{Figures 6.1--6.2}{Structure des applications Django}
\textbf{Figure 6.1}~: arborescence des 5 apps Django~: \texttt{users/} (models,
serializers, views, permissions, urls), \texttt{patients/} (models 160+ champs,
views 3800+ lignes, tasks Celery, preprocess\_rag), \texttt{predictions/}
(views, predict\_mortalite, models/feature\_mapping, train\_from\_excel),
\texttt{audit/} (models, views), \texttt{config/} (settings, urls, celery).

\textbf{Figure 6.2}~: chaque app suit le patron DRF~: Modèle (ORM Django)
$\rightarrow$ Sérialiseur (validation + transformation JSON) $\rightarrow$
Vue (logique métier APIView) $\rightarrow$ URL (routage). La séparation garantit
que modifier la validation d'un champ (sérialiseur) n'affecte pas la logique
métier (vue), et vice versa.
\end{figbox}

\begin{figbox}{Figure 6.3}{Diagramme de classes UML --- modèles Django}
Classes représentées avec leurs attributs réels (types vérifiés dans le code)~:
\begin{itemize}[leftmargin=1.5em, itemsep=1pt]
  \item \texttt{Role}~: nom~=~CharField(max=50, choices, unique=True)
  \item \texttt{Utilisateur}~: email(unique), nom, prenom, telephone CharField~;
    personal\_notes TextField~; role FK~; hérite AbstractBaseUser
  \item \texttt{Patient}~: icc\_charlson~=~\textbf{CharField} (pas Integer)~;
    utilisateur\_saisie~=~\textbf{CharField} (pas FK)~; save() auto-assign
    id\_patient~; \texttt{extract\_features\_for\_patient()} est dans
    \texttt{predictions/models/feature\_mapping.py} (pas une méthode Patient)
  \item \texttt{AuditLog}~: action~=~\textbf{TextField} (pas CharField)
  \item Relation Patient-Utilisateur~: \textbf{association simple} via
    CharField (pas agrégation, car la référence est un label texte figé)
\end{itemize}
\end{figbox}

\begin{figbox}{Tableau 6.1}{Endpoints REST principaux}
40+ endpoints groupés par ressource. Les plus utilisés dans le projet~:
\begin{center}
\small
\begin{tabular}{lll}
\toprule
\textbf{Méthode + URL} & \textbf{Vue} & \textbf{Droits} \\
\midrule
POST /api/auth/login/ & LoginView & Public \\
POST /api/auth/token/refresh/ & TokenRefreshView & Public \\
POST /api/patients/preprocess/analyze/ & PatientPreprocessAnalyzeView & Auth \\
GET /api/patients/preprocess/\{id\}/rows/ & PatientPreprocessRowsView & Auth \\
POST /api/patients/preprocess/\{id\}/integrate/ & PatientPreprocessIntegrateView & Chef \\
GET /api/predictions/patient/\{id\}/mortalite/ & PredictMortalitePatientView & Auth \\
\bottomrule
\end{tabular}
\end{center}
\end{figbox}

\begin{figbox}{Figure 6.4}{Flux JWT --- authentification et rafraîchissement}
Diagramme de séquence à 3 acteurs (Browser, Django API, PostgreSQL) en 5 étapes~:
\begin{enumerate}[leftmargin=2em, itemsep=1pt]
  \item \texttt{POST /api/auth/login/} body~: \{email, password\}
  \item \texttt{SELECT users\_utilisateur WHERE email=?} $	o$ hachage bcrypt vérifié
  \item Émission~: access JWT (8h, payload custom) + refresh JWT (1j)
  \item Requêtes authentifiées~: header \texttt{Authorization: Bearer <access>}~;
    verification HMAC-SHA256 côté Django
  \item Refresh~: \texttt{POST /api/auth/token/refresh/} + refresh token
    $	o$ nouveau access JWT (8h)
\end{enumerate}
\textbf{Important~:} l'intercepteur Axios ajoute uniquement le Bearer token
aux requêtes sortantes~; il n'y a \textbf{pas} de refresh automatique sur
expiration (pas d'intercepteur RESPONSE sur 401 dans le code actuel).
\end{figbox}

\begin{figbox}{Tableau 6.2}{Matrice de permissions RBAC}
Croise 4 rôles $\times$ 12 actions. Exemples des règles les plus importantes~:
\textbf{Super Admin}~: seul à pouvoir supprimer des comptes Super Admin et voir
tous les AuditLogs (y compris masqués). \textbf{Chef de Service}~: peut approuver/rejeter
les imports et voir les AuditLogs mais pas ceux masqués.
\textbf{Professeur/Résident}~: peuvent importer et prédire mais voient uniquement
leur propre fil d'activité (pas les AuditLogs globaux). Implémentée dans
\texttt{users/permissions.py}~: classes \texttt{IsSuperAdmin} et
\texttt{IsAdminOrChefService}, appliquées via
\texttt{permission\_classes = [IsAdminOrChefService]} dans chaque vue.
\end{figbox}

\begin{figbox}{Figures 6.5 à 6.36}{Captures d'écran de l'interface (30 figures)}
\textbf{Landing page (6.5)}~: présente les 3 modules (gestion patients, IA
prétraitement, prédiction ML), accessible sans authentification.

\textbf{Login (6.6)}~: split-panel. Côté droit~: formulaire JWT.
Après connexion réussie~: \texttt{jwtDecode(access)} extrait id, email, nom, rôle~;
\texttt{navigate('/dashboard')} redirige.

\textbf{Sidebar (6.7)}~: items de navigation filtrés selon \texttt{user.role}
via \texttt{ProtectedRoute}.

\textbf{Dashboard Super Admin (6.9--6.12)}~: 7 KPI cards interactives~;
création/modification/suppression de comptes avec mot de passe de confirmation
obligatoire (implémenté dans la vue DRF~: \texttt{confirm\_admin\_password(request)}).

\textbf{Prétraitement (6.18--6.25)}~: interface complète du pipeline LLM~:
upload $	o$ barre de progression temps réel (polling toutes les 2s) $	o$ tableau
de révision cellule par cellule $	o$ export Excel coloré $	o$ intégration PostgreSQL.

\textbf{Simulation clinique (6.34)}~: le clinicien modifie les valeurs des 32
features et voit instantanément la nouvelle prédiction via
\texttt{POST /api/predictions/patient/\{id\}/mortalite/} (mode simulation).

\textbf{Résultat prédiction (6.35)}~: jauge semi-circulaire (zone + $\hat{p}$),
top-5 facteurs (coefficients SVM), texte de recommandation personnalisé avec
taux de mortalité observé dans la zone correspondante sur HD-478.

\textbf{Monitor (6.36)}~: chronologie longitudinale par patient~: évolution des
prédictions, modifications de champs, actions des soignants --- alimenté par
les entrées AuditLog.
\end{figbox}

\begin{figbox}{Figures 6.37 à 6.48}{Graphiques analytiques et diagrammes de séquence}
\textbf{Graphiques analytiques (6.37--6.48, Chart.js)~:}
\begin{itemize}[leftmargin=1.5em, itemsep=1pt]
  \item \textbf{6.37 --- Tendance mensuelle~:} nb inclusions par mois 2020--2024,
    révèle les pics saisonniers et la croissance de la cohorte.
  \item \textbf{6.38 --- Sexe + KPIs~:} donut Homme/Femme + 4 cartes~:
    âge médian, albumine moyenne, durée dialyse médiane, taux mortalité 1 an.
  \item \textbf{6.43 --- Zones de risque~:} camembert Faible/Modéré/Élevé
    avec taux observés (3{,}8\% / 15{,}6\% / 46{,}5\%).
  \item \textbf{6.45 --- Kaplan-Meier~:} 3 courbes de survie par zone,
    log-rank $p<0{,}001$.
  \item \textbf{6.46 --- Facteurs contributifs~:} top-10 coefficients SVM
    ordonnés par valeur absolue --- base de l'interprétabilité clinique.
  \item \textbf{6.47 --- Radar profil de risque~:} profil patient (rouge) vs.\
    médiane cohorte (bleu) sur 8 dimensions.
\end{itemize}

\textbf{Diagrammes de séquence (6.48--6.50)~:}
Voir les figures détaillées dans les sections Chapitres précédentes.
Les 3 séquences couvrent~: authentification JWT (5 étapes, 3 acteurs),
import LLM (18 étapes, 7 acteurs, Celery asynchrone), prédiction SVM
(11 étapes, 5 acteurs, 1 seul appel predict\_proba).
\end{figbox}

\bigskip\noindent\rule{\linewidth}{0.4pt}\bigskip

\noindent\textbf{Résumé~:} Ce guide documente \textbf{50+ figures} et
\textbf{14 outils/technologies} avec leurs rôles précis dans le projet.
Chaque valeur numérique citée (AUC, seuils, paramètres, ports) est conforme
au code source du projet. Chaque référence bibliographique a été vérifiée
sur PubMed.

\end{document}
